\documentclass{article}
\usepackage{graphicx}   % para figuras
\usepackage{amsmath}    % para equações
\usepackage{float}      % para posicionamento de figuras/tabelas



\begin{document}

Exercício Dia 2: Animais e Energia
Kawã Mendes
18 de Outubro de 2025

Seção - Introdução
Neste exercício, vamos explorar como diferentes animais consomem energia e como podemos representar isso com figuras, tabelas e equações.  
Você vai praticar a organização do documento com seções e subseções, inserção de figuras, tabelas e uma equação simples.

subseção - Objetivos
    Criar seções e subseções.
    Inserir figuras e tabelas no documento.
    Escrever uma equação simples no meio do texto.


Subseção - Consumo de Energia em Animais
O consumo de energia varia bastante entre diferentes espécies. Animais maiores tendem a consumir mais energia absoluta, mas menos energia relativa por quilo de peso.

Subseção - Exemplo de Figura
A Figura abaixo mostra três animais diferentes: um coelho, um cachorro e um elefante.

%<COLOQUE AQUI A SUA FIGURA>

Subseção - Equação de Energia
A energia total consumida por um animal pode ser aproximada pela equação:

    E = m . c

onde E é a energia consumida, m é a massa do animal e c é o consumo médio por quilo.

Seção - Tabela de Consumo
A Tabela abaixo apresenta valores aproximados de consumo energético diário para alguns animais:

        Animal  Massa (kg)  Consumo Diário (kcal) 
        Coelho  2  200 
        Cachorro  15  1200 
        Elefante  5000  500000 

Seção - Conclusão
Neste exercício, aprendemos a organizar um documento usando seções e subseções, inserir figuras e tabelas, e colocar uma equação simples. 

\end{document}
